%% Master CV - Siddhant Gupta
%% Comprehensive overview of all competencies, experience, and achievements.
%% Use as a master reference when tailoring role-specific CVs.
%%
%% SETUP: Populated by /setup
%% Compile with: lualatex main_example.tex
%% (pdflatex often fails on modern MiKTeX with fontawesome5 font-expansion errors.)

\documentclass[11pt,a4paper,sans]{moderncv}
\moderncvstyle{banking}
\moderncvcolor{blue}

\usepackage[scale=0.80]{geometry}
\usepackage{needspace} % Added to dynamically manage awkward page breaks

% personal data
\name{Siddhant}{Gupta}
\address{Trondheimsveien 64, Oslo, Norway}{}{}
\phone[mobile]{+47 92206297}
\email{siddhantbenz@gmail.com}
\social[linkedin]{siddhantgupta1996}

\begin{document}

\makecvtitle

% ============================================================
%     PROFILE STATEMENT
% ============================================================

\vspace{6pt}
{\small Data analyst with 4+ years of experience in greenhouse gas accounting, ESG reporting, and environmental data analysis. Builds automated pipelines and ML models that turn complex datasets into decision-ready insights aligned with CSRD, GHG Protocol, and TCFD frameworks. Specialises in helping organisations measure, report, and reduce their climate footprint.}

% ============================================================
%     CORE COMPETENCIES
% ============================================================

\section{Core Competencies}
\vspace{1pt}
\begin{itemize}

\item \textbf{Data Engineering \& Analysis}: Python (pandas, NumPy, scikit-learn), SQL, ETL pipeline development, data modelling, statistical analysis.

\item \textbf{ML \& NLP}: Transformer-based NLP models for text classification, automated GHG Scope assignment, model deployment and integration.

\item \textbf{ESG \& Carbon Accounting}: GHG Protocol (Scope 1, 2, 3), CSRD/ESRS, VSME, GRI, TCFD, ISO 14064, lifecycle assessment (LCA).

\item \textbf{Visualisation \& BI}: PowerBI, Microsoft Excel, interactive dashboard development, technical documentation, stakeholder reporting.

\item \textbf{Energy Systems Modelling}: Geospatial analysis (OpenStreetMap, GIS), optimisation modelling (GAMS), simulation modelling, renewable energy potential assessment.

\end{itemize}

% ============================================================
%     PROFESSIONAL EXPERIENCE
% ============================================================

\section{Professional Experience}
\vspace{3pt}

% --- Most Recent Role ---
\cventry{05/2022--Present}{Climate Footprint Analyst}{Ducky AS}{Norway}{}{
\begin{itemize}
    \item Engineered ETL pipelines to integrate ERP system data with emission factor databases, building automated greenhouse gas (GHG) accounting aligned with VSME and ESRS frameworks
    \item Developed NLP-based machine learning models to automatically assign GHG Protocol Scope and emission factors to financial transactions, improving the accuracy and speed of greenhouse gas calculations
    \item Processed and analysed regional telecommunication data to model local commuting patterns, delivering actionable data insights to support sustainable urban mobility initiatives across 5+ municipalities
    \item Developed a lifecycle assessment model (LCA) to compare the environmental impact of various equipment ownership scenarios (reuse, rental, and buying new), demonstrating a 25--30\% reduction in lifecycle greenhouse gas emissions to support circular business models
    \item Authored technical and methodological documentation for 3 core products, translating complex data methodologies for technical and non-technical stakeholders
\end{itemize}}

\vspace{3pt}

% --- Previous Role ---
\cventry{02/2021--09/2021}{Research Assistant}{greenventory GmbH / Fraunhofer ISE}{Freiburg, Germany}{}{
\begin{itemize}
    \item Built a geospatial data model using Python and OpenStreetMap to identify optimal installation sites for heat pump integration
    \item Programmed a simulation model using Python to estimate residential geothermal heating potential, providing quantitative data to support municipal heating transition strategies
    \item Developed interactive PowerBI dashboards consolidating ESG and economic KPIs to help municipal planners make data-driven decisions
\end{itemize}}

\vspace{3pt}

% --- Earlier Role ---
% Dynamically check if there is enough space left on page 1 for the whole entry
\needspace{6\baselineskip}
\cventry{01/2019--07/2019}{Energy Analyst Intern}{Centre for Science and Environment}{New Delhi, India}{}{
\begin{itemize}
    \item Analysed operational and emissions datasets covering 70+ thermal power plants in India to identify environmental compliance trends
    \item Translated quantitative findings into policy briefs and risk assessments, outlining recommendations for pollution control technology adoption in the energy sector
\end{itemize}}

% ============================================================
%     EDUCATION
% ============================================================

\section{Education}
\vspace{1pt}

\cventry{2019--2021}{M.Sc. in Sustainable Energy Systems}{Chalmers University of Technology}{Gothenburg, Sweden}{}{
Thesis: ``Estimating the technical shallow geothermal potential for heating: A case study of L\"{o}rrach, Germany.'' Developed a geospatial simulation model to assess residential geothermal heating potential.
}

\vspace{3pt}

\cventry{2014--2018}{B.Tech. in Mechanical Engineering}{Manipal Institute of Technology}{Manipal, India}{}{
}

% ============================================================
%     LANGUAGES
% ============================================================

\section{Languages}
\vspace{1pt}
\begin{itemize}
\item English (Native/Bilingual), Hindi (Native/Bilingual), Norwegian (Limited Working), German (Limited Working), Swedish (Elementary).
\end{itemize}

% ============================================================
%     CERTIFICATIONS
% ============================================================

\section{Certifications}
\vspace{1pt}
\begin{itemize}
\item Sustainability Leadership Training (2019)
\item Urban Mobility Concepts
\item Basic Hydraulic Systems
\end{itemize}

% ============================================================
%     REFERENCES
% ============================================================

\section{References}
\vspace{1pt}
\begin{itemize}
\item{Prof. Erik Ahlgren, Department of Space, Earth and Environment, Chalmers University of Technology}
\item{Dr. Ramananda Bhatt, Professor, Manipal Institute of Technology}
\item{Karen Dundas, Senior Software Developer, Ducky AS}
\end{itemize}

\end{document}